\documentclass[12pt]{article}
\usepackage[utf8]{inputenc}
\usepackage[T1]{fontenc}
\usepackage{hyperref}
\title{Kisan Language Resources for K-12 Curriculum Development}
\author{Compiled by Claude Sonnet 5.5 (claude-sonnet-5-5) for Andr\'as Kornai}
\date{}
\begin{document}
\maketitle

\section{Introduction}
Kisan (ISO 639-3 xis) is a North Dravidian language of eastern India, mainly in Odisha, especially Sambalpur, Sundargarh and Jharsuguda. It is closely related to Kurukh (Oraon), and some sources treat it as a variety of Kurukh. Alternate names listed include Nagesia and Nagesar. A figure of 206,100 speakers (2011) is reported for the Kisan dialect. The Kisan community numbers about 331,589 people in the 2011 Scheduled Tribe census. This is a count of people, not speakers. Wikitongues lists an EGIDS status of 5 (Developing). The search found no Kisan-specific school materials.

\section{Primary Educational Resources}
The search found no Kisan-specific primers or textbooks. Because of its closeness to Kurukh, Kurukh materials may be a starting point. The search results did not show any such materials, so this is only a suggestion.

\section{Linguistic and Reference Materials}
\begin{itemize}
\item Wikitongues language page for Kisan: \url{https://wikitongues.org/languages/xis/}
\item Wikipedia, ``Kurukh language'' (covers Kisan as a related variety): \url{https://en.wikipedia.org/wiki/Kurukh_language}
\end{itemize}

\section{Religious and Community Resources}
\begin{itemize}
\item Tribes of Odisha, Kisan community profile (Odisha State Tribal Museum): \url{https://ostm.in/tribes_of_odisha/kisan/}
\end{itemize}

\section{Digital and Audio Resources}
The search found no dedicated audio or corpus resources for Kisan. The Wikitongues page may host or point to recordings, but this was not confirmed.

\section{Curriculum Development Context}
The documentation found is thin and mostly descriptive. A curriculum would need to decide whether to teach Kisan as a separate language or alongside Kurukh. It should be developed with Kisan speakers, with materials and teachers coming from the community.
\end{document}